\documentclass[11pt]{article}
\usepackage[utf8]{inputenc}
\usepackage[T1]{fontenc}
\usepackage[a4paper,margin=1in]{geometry}
\usepackage{hyperref}
\usepackage{enumitem}
\usepackage{parskip}
\setlist{nosep}

\title{Bertille Daran Website Editing Guide}
\author{}
\date{\today}

\begin{document}
\maketitle

This website is built with Hugo and the PaperMod theme.

\section*{1. Important files}
\begin{itemize}
\item \texttt{config.yml}: homepage text, menu, social links, site-wide settings
\item \texttt{static/profile.jpg}: profile photo shown on the homepage
\item \texttt{static/cv.pdf}: PDF used by the CV page
\item \texttt{static/jmp.pdf}: job market paper PDF
\item \texttt{content/research/}: research pages
\item \texttt{content/teaching/}: teaching pages
\item \texttt{content/talks/}: talks and conference pages
\item \texttt{content/location.md}: office and map page
\item \texttt{content/cv.md}: CV page with embedded PDF
\end{itemize}

\section*{2. How to change the homepage text}
Edit \texttt{config.yml}.

The main text on the homepage is in:
\begin{itemize}
\item \texttt{params.profileMode.title}
\item \texttt{params.profileMode.subtitle}
\item \texttt{params.profileMode.buttons}
\item \texttt{params.socialIcons}
\end{itemize}

Examples:
\begin{itemize}
\item change your displayed name: edit \texttt{params.profileMode.title}
\item change your bio: edit \texttt{params.profileMode.subtitle}
\item change the button labels or destinations: edit \texttt{params.profileMode.buttons}
\item change email, GitHub, LinkedIn, Scholar, or location links: edit \texttt{params.socialIcons}
\end{itemize}

\section*{3. How to change the profile picture}
Replace \texttt{static/profile.jpg} with a new image using the same filename.

If you want a different filename, also update \texttt{params.profileMode.imageUrl} in \texttt{config.yml}.

Recommended format:
\begin{itemize}
\item JPG or PNG
\item square or close to square
\item around 600 x 600 px or larger
\end{itemize}

\section*{4. How to change the menu}
Edit the \texttt{menu.main} section in \texttt{config.yml}.

Each item has:
\begin{itemize}
\item \texttt{name}: displayed label
\item \texttt{url}: destination
\item \texttt{weight}: order in the menu
\end{itemize}

\section*{5. How to edit research content}
All research pages are in \texttt{content/research/}. Each subfolder is one page bundle and must contain an \texttt{index.md}.

Current examples:
\begin{itemize}
\item \texttt{content/research/from-farm-to-cop/index.md}
\item \texttt{content/research/obesity-reviews-2023/index.md}
\item \texttt{content/research/world-development-perspectives-2022/index.md}
\end{itemize}

To edit text:
\begin{enumerate}
\item Open the page's \texttt{index.md}.
\item Update the front matter at the top.
\item Update the body text below the second \texttt{---}.
\end{enumerate}

Useful front matter fields:
\begin{itemize}
\item \texttt{title}
\item \texttt{date}
\item \texttt{summary}
\item \texttt{description}
\item \texttt{tags}
\item \texttt{author}
\end{itemize}

To add a new research page:
\begin{enumerate}
\item Create a new folder inside \texttt{content/research/}.
\item Add an \texttt{index.md} file.
\item Copy an existing research page and modify it.
\end{enumerate}

\section*{6. How to edit teaching content}
Teaching pages are in \texttt{content/teaching/}.

Current examples:
\begin{itemize}
\item \texttt{content/teaching/microeconomics-1/index.md}
\item \texttt{content/teaching/microeconomics-2/index.md}
\end{itemize}

Use the same workflow as for research pages.

\section*{7. How to edit talks and conference content}
Talk pages are in \texttt{content/talks/}.

Current examples:
\begin{itemize}
\item \texttt{content/talks/year-2026/index.md}
\item \texttt{content/talks/year-2025/index.md}
\item \texttt{content/talks/year-2024/index.md}
\item \texttt{content/talks/year-2023/index.md}
\end{itemize}

You can either keep one page per year or create one page per event.

\section*{8. How to change links to PDFs}
Static PDFs live in \texttt{static/}. Examples:
\begin{itemize}
\item \texttt{static/cv.pdf}
\item \texttt{static/jmp.pdf}
\end{itemize}

Links to those files can be written directly as \texttt{/cv.pdf} and \texttt{/jmp.pdf}.

\section*{9. How to edit the location page}
Edit \texttt{content/location.md}.

That page already contains office addresses and embedded Google Maps iframes. You can change the office name, room number, postal address, or iframe embed URL.

\section*{10. How to preview locally}
Make sure Hugo Extended is installed and available on your PATH.

Useful commands:
\begin{verbatim}
hugo server -D
hugo --gc --minify
\end{verbatim}

The preview server usually runs at \texttt{http://localhost:1313/}. Generated production files appear in \texttt{public/}.

\section*{11. How deployment works}
Deployment is handled by GitHub Actions with \texttt{.github/workflows/hugo.yml}.

Required GitHub setting:
\begin{itemize}
\item in the repository's Pages settings, the source should be \texttt{GitHub Actions}
\end{itemize}

\section*{12. Common update recipes}
\subsection*{Change your bio text}
Edit \texttt{config.yml}, then change \texttt{params.profileMode.subtitle}.

\subsection*{Change your email}
Edit the \texttt{Email} item in \texttt{params.socialIcons} inside \texttt{config.yml}.

\subsection*{Add a new paper}
\begin{enumerate}
\item Create \texttt{content/research/new-paper-slug/}
\item Add \texttt{index.md}
\item Copy front matter from an existing research page
\item Add text and links
\end{enumerate}

\subsection*{Replace the CV}
\begin{enumerate}
\item Replace \texttt{static/cv.pdf}
\item Keep the same filename if possible
\item Check \texttt{content/cv.md}
\end{enumerate}

\subsection*{Replace the profile photo}
\begin{enumerate}
\item Replace \texttt{static/profile.jpg}
\item Keep the same filename if possible
\item Check \texttt{config.yml}
\end{enumerate}

\section*{13. Notes on old demo content}
The original demo sections from the template were kept in the repo but marked as drafts so they do not surface on the rebuilt site:
\begin{itemize}
\item \texttt{content/books/}
\item \texttt{content/courses/}
\item \texttt{content/data/}
\item \texttt{content/papers/}
\end{itemize}

They can be deleted later if you no longer need them.

\end{document}
